\section*{Chapitre 2 : Contexte institutionnel --- Le ministère de la Justice du Maroc}
\addcontentsline{toc}{section}{Chapitre 2 : Contexte institutionnel --- Le ministère de la Justice du Maroc}

\bigskip

Ce chapitre ancre notre recherche dans son terrain institutionnel. Comprendre le ministère de la Justice --- ses missions, son architecture budgétaire, ses pratiques de gestion financière --- est une condition préalable indispensable à la conception d'un dispositif de prévision pertinent et opérationnellement utile. Ce chapitre s'articule autour de trois axes : une présentation de l'institution et de ses composantes, une analyse du processus budgétaire tel qu'il se déroule concrètement au sein du ministère, et enfin le cadrage du projet de recherche issu du contexte du stage.

%================================================================
\subsection*{2.1 \quad Présentation du ministère de la Justice}
\addcontentsline{toc}{subsection}{2.1 Présentation du ministère de la Justice}
%================================================================

%----------------------------------------------------------------
\subsubsection*{2.1.1 \quad Missions et organisation générale}
%----------------------------------------------------------------

Le ministère de la Justice occupe une place singulière dans l'architecture institutionnelle du Maroc. Héritier d'une tradition juridique plurielle --- islamique, coutumière et civiliste --- il est chargé de garantir le bon fonctionnement du service public de la justice, de veiller à l'application des textes législatifs et réglementaires, et d'assurer la protection des droits et libertés des citoyens. À ce titre, il constitue l'un des piliers de l'État de droit consacré par la Constitution de 2011 \citep{MEF2015}.

\paragraph{Missions institutionnelles.}

Les missions du ministère de la Justice, telles que définies par les textes organiques et réaffirmées dans les documents de programmation budgétaire, s'articulent autour de cinq domaines principaux :

\begin{itemize}
    \item \textbf{L'organisation et la gestion du système judiciaire} : supervision du réseau des juridictions sur l'ensemble du territoire national, administration du corps de la magistrature, coordination avec les autres acteurs de la chaîne pénale et civile.

    \item \textbf{La gestion du système pénitentiaire} : encadrement et supervision de la Direction Générale de l'Administration Pénitentiaire et de la Réinsertion (DGAPR), responsable des établissements pénitentiaires et des programmes de réhabilitation des détenus.

    \item \textbf{La régulation des professions juridiques} : tutelle sur les professions d'officiers ministériels --- notaires, adoul, huissiers de justice, experts judiciaires, traducteurs --- et définition des conditions d'exercice et de déontologie.

    \item \textbf{La production normative} : contribution à l'élaboration des projets de lois et de textes réglementaires, coordination des travaux législatifs en matière civile, pénale et commerciale, codification du droit marocain.

    \item \textbf{La modernisation du service public judiciaire} : développement des services numériques (e-justice), amélioration de l'accès à la justice, déploiement des programmes de réhabilitation des juridictions et de réduction des délais de traitement \citep{Ouajdouni2020}.
\end{itemize}

Ces missions s'inscrivent dans le cadre de la politique nationale de réforme de la justice, articulée autour de la Charte de la Réforme du Système Judiciaire adoptée en 2013, dont les axes ont été progressivement déclinés en programmes et actions budgétaires au titre de la LOF 130-13 \citep{Hmioui2021}.

\paragraph{Organisation générale.}

Sur le plan organisationnel, le ministère de la Justice s'articule autour de deux composantes principales : les services centraux, établis à Rabat, et les services déconcentrés répartis sur l'ensemble du territoire national.

Les \textbf{services centraux} comprennent, outre le cabinet ministériel et le secrétariat général, plusieurs directions thématiques :

\begin{itemize}
    \item La Direction des Affaires Civiles
    \item La Direction des Affaires Pénales et Criminelles
    \item La Direction des Études, de la Législation et de la Codification (DELC)
    \item La Direction des Ressources Humaines (DRH)
    \item La Direction des Affaires Financières et du Budget (DAFB)
    \item La Direction de la Réforme Numérique et de la Modernisation (DRNM)
    \item La Direction Générale de l'Administration Pénitentiaire et de la Réinsertion (DGAPR)
    \item L'Inspection Générale des Services Judiciaires (IGSJ)
\end{itemize}

Les \textbf{services déconcentrés} constituent la composante opérationnelle de l'institution. Ils comprennent :

\begin{itemize}
    \item La Cour de cassation (siège à Rabat), juridiction suprême de l'ordre judiciaire ordinaire
    \item Les Cours d'appel et les Cours administratives d'appel, réparties dans les grandes villes du Royaume
    \item Les Tribunaux de première instance (TPI) et les Sections de justice (justice de proximité)
    \item Les Tribunaux administratifs et les Tribunaux de commerce
    \item Les quelque quatre-vingts établissements pénitentiaires sous la tutelle de la DGAPR
\end{itemize}

L'ensemble du réseau judiciaire couvre l'intégralité des douze régions du Royaume, avec des implantations dans la quasi-totalité des provinces et préfectures. Cette présence territoriale étendue est l'un des facteurs qui expliquent la complexité logistique et budgétaire de l'institution.

\paragraph{Envergure budgétaire et structure des programmes.}

Par son volume de crédits et la diversité de ses programmes, le ministère de la Justice figure parmi les départements ministériels à budget significatif au sein du paysage budgétaire marocain. Conformément à la LOF 130-13, ses crédits sont structurés en plusieurs programmes correspondant chacun à une politique publique distincte \citep{MEF2015} :

\begin{itemize}
    \item \textbf{Programme 1 : Gouvernance et appui} --- regroupe les fonctions transversales de pilotage, de gestion administrative et d'encadrement institutionnel (secrétariat général, ressources humaines, systèmes d'information, communication).

    \item \textbf{Programme 2 : Services judiciaires} --- finance le fonctionnement des juridictions et les activités des magistrats, greffiers et auxiliaires de justice, ainsi que les programmes d'accès à la justice.

    \item \textbf{Programme 3 : Administration pénitentiaire et réinsertion} --- couvre les dépenses de la DGAPR (alimentation et santé des détenus, personnel pénitentiaire, sécurité des établissements, programmes de réinsertion).

    \item \textbf{Programme 4 : Réhabilitation des juridictions et des établissements pénitentiaires} --- regroupe les dépenses d'investissement en infrastructure judiciaire et pénitentiaire (construction, réhabilitation, équipement).
\end{itemize}

La masse salariale --- magistrats, greffiers, personnel pénitentiaire --- représente la composante la plus importante et la plus rigide du budget, concentrant historiquement plus de 70\% des crédits de fonctionnement \citep{MJMaroc2023}. Les dépenses d'investissement, quant à elles, se caractérisent par des profils d'exécution particulièrement irréguliers, avec des taux de paiement concentrés sur le dernier trimestre de l'exercice --- phénomène constitutif du problème de prévision que cette étude cherche à résoudre.

\begin{center}
\begin{tabular}{llr}
\toprule
\textbf{Programme} & \textbf{Intitulé} & \textbf{Poids approximatif} \\
\midrule
Programme 1 & Gouvernance et appui & $\sim$\,8\,\% \\
Programme 2 & Services judiciaires & $\sim$\,22\,\% \\
Programme 3 & Administration pénitentiaire et réinsertion & $\sim$\,55\,\% \\
Programme 4 & Réhabilitation des infrastructures & $\sim$\,15\,\% \\
\bottomrule
\end{tabular}
\end{center}
\vspace{4pt}
{\small\textit{Tableau 2.1 --- Structure des programmes budgétaires du ministère de la Justice (part approximative en crédits ouverts)}}

%----------------------------------------------------------------
\subsubsection*{2.1.2 \quad La Direction des Affaires Financières et du Budget : rôle et positionnement dans la chaîne de dépense}
%----------------------------------------------------------------

Au sein de l'organigramme ministériel, la Direction des Affaires Financières et du Budget (DAFB) joue un rôle central dans la gestion des ressources financières de l'institution. C'est elle qui constitue l'interface principale entre le ministère de la Justice d'une part, et le ministère de l'Économie et des Finances (MEF) et la Trésorerie Générale du Royaume (TGR) d'autre part.

\paragraph{Positionnement institutionnel.}

La DAFB est rattachée hiérarchiquement au Secrétariat Général du ministère. Elle assure la coordination des fonctions financières entre les différentes directions centrales et les services déconcentrés, en jouant un rôle d'agrégateur de l'information budgétaire et de conseil auprès des responsables de programme. Son positionnement au carrefour de la chaîne de dépense lui confère une vision consolidée des flux financiers de l'institution, ce qui en fait l'interlocuteur naturel pour tout projet d'amélioration des outils de prévision budgétaire.

\paragraph{Fonctions principales.}

Les fonctions de la DAFB couvrent l'ensemble du cycle de vie budgétaire :

\begin{itemize}
    \item \textbf{La programmation budgétaire} : coordination de l'élaboration des projets de performance annuels et triennaux (PBT-CDMT), consolidation des besoins remontés par les centres de responsabilité, et défense des demandes de crédits lors des conférences budgétaires avec la Direction du Budget du MEF.

    \item \textbf{L'exécution budgétaire} : gestion des crédits délégués, visa et contrôle des engagements de dépenses, coordination avec les ordonnateurs secondaires déconcentrés, et suivi du taux de consommation des crédits en temps réel via le système GID de la TGR.

    \item \textbf{Le reporting et le contrôle de gestion} : production des tableaux de bord budgétaires périodiques, renseignement des indicateurs de performance contractualisés avec le MEF, et alimentation des Rapports Annuels de Performance (RAP) transmis au Parlement conformément à la LOF.

    \item \textbf{La gestion des morasses et des archives financières} : extraction, compilation et archivage des états d'exécution budgétaire mensuels (les \textit{morasses}), qui constituent la matière première de notre analyse empirique.

    \item \textbf{La gestion des restes à payer (RAP)} : suivi des engagements pluriannuels non encore soldés, coordination avec la TGR pour les ordonnancements différés, et pilotage de la trésorerie budgétaire.
\end{itemize}

\paragraph{Interfaces externes.}

Dans l'exercice de ses missions, la DAFB interagit avec plusieurs acteurs extérieurs au ministère. La \textbf{Direction du Budget} du MEF fixe les enveloppes dans les lettres de cadrage et arbitre les demandes lors des conférences budgétaires annuelles. La \textbf{Trésorerie Générale du Royaume (TGR)} assure l'exécution comptable et le contrôle des dépenses publiques via le système GID, dont la DAFB est l'utilisateur principal au sein du ministère. L'\textbf{Inspection Générale des Finances (IGF)} procède à des audits périodiques de la gestion financière. La \textbf{Cour des Comptes} exerce le contrôle juridictionnel \textit{a posteriori} sur les comptes publics, et ses observations annuelles alimentent le dialogue de gestion entre le ministère et le MEF.

%================================================================
\subsection*{2.2 \quad Le processus budgétaire au sein du ministère}
\addcontentsline{toc}{subsection}{2.2 Le processus budgétaire au sein du ministère}
%================================================================

%----------------------------------------------------------------
\subsubsection*{2.2.1 \quad Cycle de préparation et d'exécution}
%----------------------------------------------------------------

Le processus budgétaire au ministère de la Justice, comme dans l'ensemble des ministères marocains depuis l'entrée en vigueur de la LOF 130-13, s'articule autour d'un cycle annuel structuré dont chaque phase mobilise différents niveaux hiérarchiques et fonctionnels.

\paragraph{Phase 1 --- La programmation stratégique (janvier--avril).}

Le cycle s'ouvre, en début d'année $N$, par une phase de programmation pluriannuelle. Chaque direction centrale du ministère et chaque responsable de programme procède à la révision de son Cadre de Dépenses à Moyen Terme (CDMT) sur l'horizon $N+1$/$N+2$/$N+3$, en actualisant les prévisions de crédits à la lumière des réalisations de l'exercice $N-1$ et des nouvelles priorités stratégiques. Cette phase mobilise des méthodes de projection tendancielle pour les dépenses de fonctionnement et des estimations de coûts de projets pour les dépenses d'investissement. C'est précisément à ce stade que les limites des approches classiques de prévision se manifestent le plus clairement, comme l'a établi l'analyse de la section 1.2.

\paragraph{Phase 2 --- La conférence budgétaire et les arbitrages (mai--juillet).}

Au mois de mai, la Direction du Budget du MEF adresse aux ministères sectoriels une circulaire budgétaire fixant les lettres de plafond et les hypothèses macroéconomiques pour l'élaboration du Projet de Loi de Finances (PLF) de l'année $N+1$. Sur cette base, le ministère de la Justice procède à ses arbitrages internes --- entre programmes, entre fonctionnement et investissement, entre directions concurrentes --- avant de soumettre ses propositions lors des conférences budgétaires organisées par le MEF. Cette phase de négociation est cruciale pour la pertinence des prévisions pluriannuelles : les plafonds de crédits finalement retenus conditionnent la programmation sur $N+1$ et $N+2$, et une révision à la baisse des enveloppes génère mécaniquement des écarts entre prévision initiale et dotation réelle.

\paragraph{Phase 3 --- Le vote parlementaire et la promulgation (août--décembre).}

Le PLF, consolidant les propositions arbitrées de l'ensemble des ministères, est déposé au Parlement en octobre. Après discussion en commissions et vote en séance plénière par les deux Chambres, la Loi de Finances est idéalement promulguée avant la fin de l'exercice en cours pour permettre une mise en exécution dès le 1\textsuperscript{er} janvier de l'exercice suivant.

\paragraph{Phase 4 --- L'exécution budgétaire (janvier--décembre de l'année $N+1$).}

L'exécution s'ouvre par la délégation de crédits aux ordonnateurs secondaires déconcentrés (cours d'appel, DGAPR régionale, etc.). Elle suit ensuite le cycle séquentiel de la dépense --- engagement, liquidation, ordonnancement, paiement --- tel que décrit en section 1.1.3. L'exécution budgétaire au ministère de la Justice est caractérisée par une saisonnalité marquée : les dépenses de fonctionnement (masse salariale mensuelle, consommations courantes) sont relativement régulières, tandis que les dépenses d'investissement présentent un profil fortement concentré sur le dernier trimestre, particulièrement les mois de novembre et décembre, en raison des délais de passation des marchés publics et de la pression à consommer les crédits avant la clôture de l'exercice.

\paragraph{Phase 5 --- L'évaluation et le reporting (premier semestre $N+2$).}

La clôture de l'exercice donne lieu à l'élaboration du Rapport Annuel de Performance (RAP), qui compare les réalisations aux objectifs et indicateurs contractualisés. Ce document est transmis au Parlement et alimente la mise à jour des projections pluriannuelles du cycle suivant.

\begin{center}
\begin{tabular}{lll}
\toprule
\textbf{Période} & \textbf{Phase} & \textbf{Acteurs principaux} \\
\midrule
Jan.--Avr. & Programmation pluriannuelle (CDMT-PBT) & DAFB, Responsables de programme \\
Mai--Juil. & Conférence budgétaire et arbitrages & DAFB, Direction du Budget (MEF) \\
Août--Déc. & Vote parlementaire et promulgation & Parlement, MEF, DAFB \\
Jan.--Déc. & Exécution budgétaire (GID) & DAFB, ordonnateurs, TGR \\
Jan.--Juin & Bilan et rapport annuel de performance & DAFB, Secrétariat Général \\
\bottomrule
\end{tabular}
\end{center}
\vspace{4pt}
{\small\textit{Tableau 2.2 --- Cycle budgétaire annuel au ministère de la Justice}}

%----------------------------------------------------------------
\subsubsection*{2.2.2 \quad Outils et pratiques actuels}
%----------------------------------------------------------------

La gestion budgétaire au ministère de la Justice s'appuie sur un ensemble d'outils informatiques et de pratiques organisationnelles qui ont évolué significativement depuis l'entrée en vigueur de la LOF 130-13.

\paragraph{Le système GID (Gestion Intégrée de la Dépense).}

Le système GID, déployé et administré par la Trésorerie Générale du Royaume, constitue le référentiel central de l'exécution budgétaire dans l'ensemble des ministères marocains. Il couvre toutes les phases de la dépense publique --- engagement, liquidation, ordonnancement --- et fournit une traçabilité complète de chaque opération budgétaire depuis l'ouverture des crédits jusqu'au paiement par le comptable public.

Pour la DAFB, le GID est à la fois un outil d'exécution et une source de données. Les états d'exécution extraits mensuellement de ce système --- désignés dans la pratique administrative marocaine sous le terme de \textit{morasses} --- constituent les documents de référence pour le suivi de la consommation des crédits. Ces morasses fournissent, pour chaque programme et chaque ligne budgétaire, les montants cumulés des crédits ouverts, délégués, engagés, liquidés et ordonnancés à la date d'extraction.

\paragraph{Le portail e-Budget.}

Le portail e-Budget du MEF est utilisé dans la phase de programmation pluriannuelle pour la saisie et la consolidation des propositions budgétaires lors des conférences annuelles. Il centralise les informations relatives aux enveloppes CDMT, aux projets de performance et aux arbitrages entre le ministère et la Direction du Budget. Bien qu'il marque un progrès important en matière de standardisation des échanges d'information entre administrations, il présente des fonctionnalités d'analyse prospective encore limitées par rapport aux besoins opérationnels des services financiers ministériels.

\paragraph{Les outils bureautiques.}

Malgré la disponibilité de systèmes structurés comme le GID et e-Budget, une part significative du travail de programmation et d'analyse budgétaire demeure réalisée à l'aide de tableurs (Microsoft Excel). Les responsables de programme et les agents de la DAFB recourent à des fichiers Excel pour consolider les données issues du GID, calculer les indicateurs de taux de consommation et élaborer les projections pluriannuelles selon des méthodes tendancielles. Cette dépendance aux tableurs, si elle répond à un besoin réel de flexibilité analytique, génère des risques d'erreurs de saisie, de silos d'information entre services, et d'absence de traçabilité et de reproductibilité des analyses \citep{Hmioui2021}. C'est l'une des faiblesses que notre projet vise à contribuer à corriger.

\paragraph{Le reporting budgétaire.}

Le ministère produit, à intervalles réguliers, plusieurs types de documents de reporting : des tableaux de bord mensuels à usage interne (situation des engagements, taux de consommation par programme et par nature de dépense), des rapports trimestriels à destination du MEF, et des Rapports Annuels de Performance (RAP) transmis au Parlement conformément aux exigences de la LOF. La production de ces documents mobilise un travail manuel important de consolidation et de vérification des données issues du GID, ce qui illustre l'espace disponible pour des outils plus automatisés et plus robustes.

%----------------------------------------------------------------
\subsubsection*{2.2.3 \quad Diagnostic de l'existant : forces et limites}
%----------------------------------------------------------------

L'analyse du processus budgétaire du ministère de la Justice révèle un ensemble de forces et de limites qui définissent précisément le contexte opérationnel dans lequel s'inscrit notre projet de recherche.

\paragraph{Acquis institutionnels.}

Depuis la mise en œuvre progressive de la LOF 130-13 (2016--2020), le ministère de la Justice a accompli des progrès notables dans la modernisation de sa gestion budgétaire :

\begin{itemize}
    \item La structuration des crédits en programmes (services judiciaires, administration pénitentiaire, réhabilitation des infrastructures, gouvernance) a clarifié les responsabilités budgétaires et amélioré la lisibilité de l'allocation des crédits pour les parlementaires et les auditeurs externes.

    \item Les contrats de performance conclus avec le MEF ont introduit une culture d'accountability qui commence à infuser les pratiques managériales des responsables de programme.

    \item Le déploiement généralisé du GID a uniformisé les pratiques d'exécution et renforcé la traçabilité des dépenses à toutes les étapes du cycle budgétaire.

    \item La constitution progressive d'un historique de données d'exécution (plusieurs exercices de morasses mensuelles depuis 2016) crée le substrat technique indispensable à des analyses prospectives avancées --- substrat dont notre étude exploite précisément le potentiel.
\end{itemize}

\paragraph{Limites et points de tension.}

Plusieurs fragilités structurelles persistent néanmoins, qui justifient la recherche de méthodes de prévision plus performantes \citep{Abdelali2019} :

\textbf{La qualité et la précision des prévisions budgétaires} demeurent insuffisantes au regard des exigences de sincérité imposées par la LOF. Les écarts entre les prévisions initiales et les réalisations effectives sont récurrents et parfois significatifs, en particulier pour les dépenses d'investissement (programme 4) et les dépenses de fonctionnement hors masse salariale. Ces écarts alimentent un cycle problématique : sous-consommation de crédits en cours d'exercice, demandes de reports ou de virements en fin d'année, et distorsions dans la calibration des enveloppes de l'exercice suivant.

\textbf{La réactivité des outils de prévision} est limitée. Les méthodes de projection actuelles, principalement tendancielles, peinent à intégrer rapidement les nouvelles informations pertinentes --- évolution de la population carcérale, réorientations de politique pénale, chocs exogènes --- dans les projections budgétaires. La pandémie de COVID-19 de 2020-2021 a illustré de manière particulièrement saillante les limites des modèles fondés sur la simple continuité des tendances historiques.

\textbf{La fragmentation de l'information} constitue un obstacle récurrent à une prévision de qualité. Les données utiles sont dispersées entre plusieurs systèmes (GID pour l'exécution, e-Budget pour la programmation, systèmes sectoriels de la DGAPR pour les données pénitentiaires, outils de reporting du service des ressources humaines pour la masse salariale) et leur consolidation demeure largement manuelle, coûteuse en temps et source d'erreurs.

\textbf{Le manque de compétences en analyse quantitative avancée} au sein des services financiers constitue un frein à l'adoption spontanée d'outils plus sophistiqués. Si les agents maîtrisent parfaitement les outils réglementaires et comptables, les capacités en modélisation statistique et en science des données restent peu développées, ce qui renforce la nécessité d'accompagner tout outil d'aide à la décision d'une démarche de formation et de transfert de compétences.

%================================================================
\subsection*{2.3 \quad Contexte du stage et cadrage du projet}
\addcontentsline{toc}{subsection}{2.3 Contexte du stage et cadrage du projet}
%================================================================

%----------------------------------------------------------------
\subsubsection*{2.3.1 \quad Enjeux identifiés lors du stage}
%----------------------------------------------------------------

Le stage réalisé au sein de la Direction des Affaires Financières et du Budget du ministère de la Justice a constitué le point d'entrée opérationnel de ce projet de recherche. Il a permis d'observer de l'intérieur les pratiques de programmation et d'exécution budgétaire, de participer aux travaux de consolidation des morasses et d'élaboration des tableaux de bord, et d'identifier avec précision les points de tension les plus significatifs.

L'enjeu central mis en évidence est celui de la \textbf{précision de la prévision budgétaire à l'horizon pluriannuel}. Les échanges avec les responsables de la DAFB ont révélé une insatisfaction partagée quant aux méthodes actuelles de projection, notamment pour les dépenses d'investissement, les dépenses de fonctionnement hors masse salariale, et --- dans une moindre mesure --- les recettes propres (amendes judiciaires, droits de timbre). Les agents témoignent de la récurrence des écarts entre prévisions et réalisations, et du temps considérable consacré aux ajustements en cours et en fin d'exercice.

Derrière cet enjeu de précision se dessine un enjeu plus large de \textbf{gouvernance budgétaire et de sincérité}. La LOF 130-13 a fait de la sincérité des prévisions et de l'atteinte des objectifs de performance deux piliers de la réforme des finances publiques. Un ministère dont les prévisions s'écartent systématiquement des réalisations s'expose à des critiques lors des évaluations parlementaires et des audits de la Cour des Comptes. La mise en place d'outils de prévision plus robustes constitue donc une réponse à la fois technique et institutionnelle à ces pressions croissantes.

Un troisième enjeu, révélé par l'analyse de la période récente, est celui de la \textbf{résilience face aux chocs exogènes}. L'épisode de la pandémie de COVID-19 en 2020-2021 a illustré de manière particulièrement saillante les limites des modèles de prévision fondés sur la continuité des tendances passées : suspension temporaire des audiences, libérations exceptionnelles de détenus réduisant les besoins pénitentiaires, gel de certains programmes d'investissement --- autant de perturbations qu'aucune méthode tendancielle n'aurait pu anticiper, mais dont les effets auraient pu être mieux encadrés grâce à des modèles intégrant des variables de contexte exogènes.

%----------------------------------------------------------------
\subsubsection*{2.3.2 \quad Objectifs du projet de recherche}
%----------------------------------------------------------------

Sur la base des enjeux identifiés et des lacunes documentées dans la littérature, ce projet de recherche poursuit quatre objectifs complémentaires et articulés.

\textbf{Objectif 1 --- Caractériser les écarts de prévision budgétaire au ministère de la Justice.}
Il s'agit de produire une analyse descriptive et statistique rigoureuse des écarts entre prévisions initiales et réalisations effectives sur la période post-LOF, d'en dégager la structure (systématique vs aléatoire), les régularités temporelles et les patterns différenciés selon les programmes.

\textbf{Objectif 2 --- Construire et entraîner des modèles de prévision alternatifs.}
À partir des données des morasses budgétaires et de variables exogènes sélectionnées, développer et calibrer plusieurs modèles de prévision appartenant à trois familles méthodologiques : économétrique (ARIMA/SARIMA, VAR, Prophet), Machine Learning (Random Forest, XGBoost) et Deep Learning (LSTM, GRU).

\textbf{Objectif 3 --- Comparer les performances prédictives des modèles.}
Sur la base d'un protocole de validation rigoureux (validation croisée temporelle walk-forward, métriques standards MAE, RMSE, MAPE, $R^2$, test de Diebold-Mariano), évaluer la supériorité relative des approches d'IA par rapport aux méthodes économétriques de référence, et identifier les conditions dans lesquelles chaque famille de modèles est la plus performante.

\textbf{Objectif 4 --- Formuler des recommandations opérationnelles et développer un prototype.}
Capitalisant sur les résultats de la modélisation, proposer des recommandations pratiques pour l'amélioration des outils de programmation budgétaire au sein de la DAFB, et développer un prototype d'outil de visualisation (Streamlit) facilitant l'interprétation et l'appropriation des prévisions par les équipes opérationnelles.

Ces quatre objectifs s'inscrivent directement dans le prolongement des trois hypothèses de recherche formulées en section 1.4, et dont la vérification empirique constitue le cœur des chapitres 3 à 5 de ce mémoire.
